\documentclass[11pt]{article}

\usepackage[margin=1in]{geometry}
\usepackage{booktabs}
\usepackage{amsmath}
\usepackage{amssymb}
\usepackage{graphicx}
\usepackage{hyperref}
\usepackage{natbib}
\usepackage{xcolor}

\emergencystretch=3em

\title{Label-Free Event-Aware LLM Adaptation for Post-Pandemic Household Travel Demand Prediction}
\author{Anonymous Authors}
\date{}

\newcommand{\nhts}{\textsc{NHTS}}
\newcommand{\method}{\textsc{EventAdapter}}
\newcommand{\whmae}{weighted MAE}
\newcommand{\whbias}{weighted bias}

\begin{document}
\maketitle

\begin{abstract}
Household travel surveys provide calibrated evidence for transportation planning, but temporal transfer can fail when rare societal shocks change the behavioral meaning of otherwise stable household covariates.
We study this problem with the U.S. National Household Travel Survey (\nhts), training routine-mobility predictors on pre-2022 waves and evaluating on the 2022 post-pandemic recovery wave.
A historical XGBoost model substantially overpredicts 2022 household trip generation, with \whmae{} 4.3377 and \whbias{} +3.6052.
We propose a label-free event-aware adaptation framework in which a large language model (LLM) does not directly predict household trips.
Instead, it generates cohort-level event priors for mechanisms such as trip suppression, remote-work substitution, transit avoidance, delivery substitution, and recovery sensitivity.
These priors are distilled into a fixed no-label event adapter that corrects the historical routine predictor without using 2022 trip-count labels for training or calibration.
The primary gated adapter reduces \whmae{} to 2.5023, reduces weighted RMSE by 32.22\%, and nearly removes systematic overprediction with \whbias{} -0.0230.
We evaluate the method against historical, trend, global-event, random-prior, LLM-only, zero-shot rule-tree, small-history rule-calibration, and stronger non-LLM baselines.
The results suggest that LLMs are most useful in this survey setting as constrained event-prior generators rather than standalone numerical predictors.
\end{abstract}

\section{Introduction}

Household travel demand prediction is usually grounded in historical survey data.
Models trained on prior survey waves can learn routine relationships between household attributes and travel behavior: household size, vehicle ownership, worker count, income, urban form, day of week, and regional context all influence how much a household travels on a diary day.
This assumption becomes fragile when the target year is affected by a rare societal event.
The 2022 \nhts{} wave is a post-pandemic recovery wave, after remote work, public-transit avoidance, online shopping, and uneven activity recovery changed daily travel patterns \citep{fhwa2022nhts,fhwa2024nhtstrends}.

This paper asks whether an LLM can help adapt a historical household travel model to a post-pandemic target year without using target-year trip-count labels for training or calibration.
The key design choice is to separate routine mobility from event response.
A supervised tabular model estimates household-level routine demand from historical \nhts{} waves, while the LLM supplies structured event priors for how the post-pandemic context may suppress or reshape travel.
The LLM is therefore not a direct trip-count regressor.
It is a constrained event-generalization module whose outputs are validated and distilled into an auditable correction rule.

Recent 2025--2026 mobility work motivates this role separation.
ELLMob frames event-driven mobility as a conflict between habitual patterns and event constraints \citep{wang2026ellmob}.
CausalMob transforms public-event text into LLM-derived intention features for mobility prediction with a causal framing \citep{yang2025causalmob}.
AgentMove studies zero-shot mobility prediction with an LLM agentic framework \citep{feng2025agentmove}, while AgentMob points to a fast-path plus selective-reasoning pattern for efficient evidence-grounded mobility prediction \citep{chen2026agentmob}.
ELP-Mob and UniMob reflect related trends in efficient LLM mobility pipelines and universal mobility modeling \citep{wang2025elpmob,long2025unimob}.
Broader surveys show that transportation and spatio-temporal modeling are moving toward language-model and foundation-model paradigms \citep{yan2025transportllm,liang2025stfmsurvey}.
Recent supervised daily-travel work also combines recurrent neural networks and LLMs for travel behavior prediction, but it does not target the no-target-year-label adaptation setting studied here \citep{sharifi2026dailytravel}.
Our target differs from these settings: the prediction unit is a household survey record, not a GPS trajectory, road sensor, or flow node, and the target-year labels should be unavailable in the intended prospective setting.

\paragraph{Contributions.}
We make four contributions.
First, we formulate 2022 \nhts{} household trip generation as event-driven temporal adaptation rather than ordinary cross-year forecasting.
Second, we introduce a label-free LLM event-prior adapter that corrects a historical household predictor without using 2022 \texttt{CNTTDHH} labels for calibration.
Third, we evaluate against a broad method spectrum, including historical mean, ordinary XGBoost, trend shift, CatBoost GPU, LLM-only pressure, global event prior, random prior, zero-shot LLM rule tree, pseudo-label tree, and LLM rule plus historical calibration.
Fourth, we add leakage audits, placebo controls, bootstrap confidence intervals, equity-aware subgroup analysis, mode/purpose extensions, and external mechanism checks with ACS and PSRC data \citep{burrows2024commuting,psrc2024regionaltravelstudy}.

\section{Problem Definition}

Let $X$ denote harmonized household features and $y$ denote household diary-day trip count \texttt{CNTTDHH}.
Historical \nhts{} waves are available for training, and the target wave is 2022.
The strict setting forbids using 2022 \texttt{CNTTDHH} labels for training, calibration, or parameter selection.
The labels are used only for final evaluation.

The main task is:
\begin{quote}
Given historical labeled \nhts{} waves and unlabeled 2022 household profiles, predict 2022 household travel behavior under post-pandemic event shift.
\end{quote}

The primary target is household trip generation.
Additional behavior outputs include household mode-share vector, household purpose-share vector, and derived mode-specific trip volume:
\begin{equation}
\widehat{y}_{m} = \widehat{y}_{\mathrm{trips}} \cdot \widehat{s}_{m},
\end{equation}
where $\widehat{s}_{m}$ is the predicted household share for mode $m$.

\section{Method}

\subsection{Historical Routine Predictor}

We train a historical routine-mobility model $f_{\mathrm{hist}}(X)$ on pre-2022 \nhts{} household records.
The final ordinary supervised baseline is a CUDA XGBoost model trained on 2001, 2009, and 2017 and evaluated on 2022.
This branch captures household heterogeneity but lacks post-pandemic event semantics.

\subsection{LLM Event-Prior Generation}

The LLM receives cohort-level household profiles rather than individual household targets.
Cohorts summarize safe household attributes and exclude \texttt{HOUSEID}, \texttt{CNTTDHH}, survey weights, and target-derived aggregate statistics.
The output schema contains numeric event priors:
\begin{itemize}
    \item trip suppression,
    \item remote-work substitution,
    \item transit avoidance,
    \item delivery substitution,
    \item post-pandemic recovery sensitivity,
    \item confidence.
\end{itemize}
The 2022 households are compressed from 7,893 rows to 1,327 cohorts.
Batch prompting with batch size 15 reduces event-prior generation to 89 prompts.
After schema and range validation, inference uses a deterministic adapter rather than per-household online LLM calls.

\subsection{Fixed No-Label Adapter}

The primary prediction is
\begin{equation}
\widehat{y}_{2022}
= f_{\mathrm{hist}}(X)
\cdot
\mathrm{clip}\left(1 - \alpha s_{\mathrm{event}}, \lambda_{\min}, 1\right),
\label{eq:adapter}
\end{equation}
where $s_{\mathrm{event}}$ is the distilled event pressure.
The reported primary rule is the gated trip-suppression adapter.
Low-confidence or weakly differentiated cohort priors fall back to a global event pressure.
This keeps the main experiment fixed and label-free.

\subsection{What the Method Is Not}

The method is not pure LLM regression, not an LLM-generated 2022 decision tree, and not a generic mobility foundation model.
It is a hybrid structured-data system: historical labels provide the numerical household baseline, and LLM event priors provide semantic adaptation under post-pandemic shift.

\section{Experimental Setup}

\paragraph{Data.}
The main dataset is public-use \nhts{} household data for 2001, 2009, 2017, and 2022, with 357,553 harmonized household records and 28 common household variables.
The target 2022 evaluation set contains 7,893 households.
External mechanism checks use ACS commuting statistics and PSRC household travel survey data \citep{burrows2024commuting,psrc2024regionaltravelstudy}.

\paragraph{Metrics.}
For trip generation, we report weighted MAE, weighted RMSE, weighted bias, weighted $R^2$, and household-level tolerance accuracy.
Weighted bias is central because historical transfer overpredicts total demand.
For mode and purpose composition, we report weighted total variation and component MAE.

\paragraph{Baselines and ablations.}
The comparison includes historical mean, ordinary historical XGBoost, historical trend shift, stronger tabular baselines, LLM-only pressure, global event prior, random pressure control, zero-shot LLM rule tree, pseudo-label tree, LLM rule plus historical calibration, and the primary gated LLM event adapter.

\section{Results}

\subsection{Main Trip-Generation Result}

\begin{table}[t]
\centering
\caption{Household trip-generation results on the 2022 \nhts{} target wave. Lower \whmae{}, RMSE, and absolute bias are better.}
\label{tab:main_results}
\begin{tabular}{lrrrr}
\toprule
Method & \whmae{} & weighted RMSE & \whbias{} & weighted $R^2$ \\
\midrule
Historical mean & 5.5116 & 6.1251 & +4.4995 & -1.1722 \\
Ordinary XGBoost & 4.3377 & 5.3169 & +3.6052 & -0.6367 \\
Historical trend shift & 4.1504 & 5.1463 & +3.3485 & -0.5334 \\
LLM-only pressure & 2.7175 & 3.9876 & -0.3732 & 0.0794 \\
Global event prior & 2.5223 & 3.7432 & -0.7477 & 0.1888 \\
Random prior control & 2.6466 & 3.8807 & -0.6368 & 0.1280 \\
Zero-shot LLM rule tree & 2.6019 & 3.8300 & -0.4072 & 0.1507 \\
LLM rule + 500-history calibration & 2.7723 & 3.8226 & +0.4627 & 0.1538 \\
\method{} gated adapter & \textbf{2.5023} & \textbf{3.6038} & \textbf{-0.0230} & \textbf{0.2480} \\
\bottomrule
\end{tabular}
\end{table}

The ordinary XGBoost baseline overpredicts strongly.
The primary gated adapter reduces \whmae{} by 42.31\%, weighted RMSE by 32.22\%, and absolute weighted bias by 99.36\%.
The LLM-only pressure baseline is directionally useful but less calibrated, supporting the hypothesis that LLM priors need a historical household baseline.

\subsection{Statistical Validation}

Household bootstrap resampling gives 95\% confidence intervals:
ordinary XGBoost \whmae{} [4.2418, 4.4312], gated adapter \whmae{} [2.4292, 2.5808], paired \whmae{} reduction [1.7422, 1.9267], and paired absolute-bias reduction [3.3582, 3.6696].
These intervals support that the improvement is not a single point-estimate artifact.

\subsection{Controls and Robustness}

The global event prior is strong, so the correct interpretation is event-level adaptation plus incremental cohort refinement, not ``LLM ranking explains everything.''
Placebo and robustness checks constrain the claim:
\begin{itemize}
    \item 500-run permutation control for the primary gated rule: actual \whmae{} 2.5023, random-permutation mean 2.5808, empirical $p=0.0020$.
    \item Irrelevant pseudo-event controls are weaker: best ranked pseudo-event \whmae{} 2.6274, best gated pseudo-event \whmae{} 2.5610.
    \item Zero-shot LLM rule tree reaches \whmae{} 2.6019, and a pseudo-label tree distilled from it reaches 2.6040; both are weaker and more biased than the primary adapter.
    \item LLM rule plus 500 historical calibration reaches \whmae{} 2.7723, showing that rule distillation plus small historical calibration is coherent but not the best solution for the 2022 event-shift task.
    \item CatBoost GPU is the strongest extra non-LLM baseline with \whmae{} 4.2196 and \whbias{} +3.4873.
\end{itemize}

\subsection{Behavior-System Outputs}

The method extends beyond total trips.
Transit-share weighted MAE improves from 0.0325 to 0.0269.
Total mode-trip MAE improves from 4.8467 to 3.2802 when combining gated trip count with LLM-adjusted mode shares.
Purpose composition is implemented as an exploratory third dimension, but the LLM prior is not the overall best purpose row.
These results support a broader household behavior-system framing, while trip generation remains the main contribution.

\subsection{External and Temporal Validation}

Temporal transfer checks show that 2022 is more difficult than pre-COVID transfer: pre-COVID mean absolute weighted bias is 0.9024, while a diagnostic 2022 transfer check has absolute weighted bias 3.7202.
External evidence supports the event mechanisms.
ACS commuting statistics show worked-from-home share 5.7\% to 15.2\% from 2019 to 2022, while public-transportation commute share changes 5.0\% to 3.1\% \citep{burrows2024commuting}.
PSRC household survey replication shows a pre/post improvement from \whmae{} 3.4271 to 3.2498 and \whbias{} +1.0532 to +0.2605 using an ACS-derived remote-work suppression factor without target-year PSRC label calibration.
The PSRC result is an external regional replication of the event-adaptation principle, not direct numerical validation of \nhts{} 2022.

\section{Discussion}

The main finding is role separation.
Historical tabular models supply calibrated household heterogeneity but fail under post-pandemic event shift.
LLM priors supply event semantics but are under-calibrated as standalone numerical predictors.
The hybrid adapter works because it constrains the LLM to a narrow semantic role and leaves numerical scale to historical survey data.

The strongest alternative explanation is that a global post-pandemic downscaling scalar captures most of the benefit.
The experiments partly confirm this: the global event prior is strong.
The contribution should therefore be framed carefully as label-free event adaptation with auditable cohort refinement, not as evidence that rich LLM reasoning alone dominates.
Placebo controls, zero-shot rule-tree baselines, small-calibration baselines, and subgroup checks provide evidence that the event prior is meaningful, but additional regional and shock replications would strengthen the paper.

\section{Limitations}

The current evidence has five limitations.
First, the main \nhts{} evaluation focuses on one shock year, 2022.
Second, the LLM may contain broad post-pandemic mobility knowledge from pretraining; the prospective event-context file reduces but does not fully eliminate retrospective knowledge concerns.
Third, the global event prior is strong, so cohort-specific LLM ranking should be described as incremental.
Fourth, mode and purpose extensions are useful for planning relevance but are not yet full mode-choice or purpose-choice models.
Fifth, PSRC is a regional external survey with different sampling and diary protocols, so it validates transfer of the principle rather than exact \nhts{} numerical accuracy.

\section{Conclusion}

This project shows that post-pandemic household travel prediction can be formulated as label-free event adaptation.
A historical household model alone overpredicts 2022 travel, while a pure LLM-style predictor is directionally useful but under-calibrated.
Distilling LLM event semantics into a fixed adapter substantially reduces trip-count error and nearly removes systematic bias without using 2022 trip-count labels for calibration.
The result suggests a practical path for using LLMs in survey-based mobility modeling: constrain them to produce auditable event priors, combine them with historically calibrated models, and evaluate them with causal guardrails, robust baselines, and planning-oriented multi-objective metrics.

\section*{Artifact Notes}

The numerical results in this draft are backed by repository artifacts:
\begin{itemize}
    \item \texttt{outputs/final\_project/final\_metrics\_summary.csv}
    \item \texttt{outputs/statistical\_validation/}
    \item \texttt{outputs/robustness\_checks/}
    \item \texttt{outputs/zero\_shot\_llm\_rule\_tree\_baseline/}
    \item \texttt{outputs/llm\_rule\_small\_data\_calibration/}
    \item \texttt{outputs/external\_validation/}
\end{itemize}

\bibliographystyle{plainnat}
\bibliography{references}

\end{document}
